\section{Experiments}
\label{sec:experiments}

\subsection{Experimental Setup}
\label{sec:exp-setup}

\paragraph{Teacher--student configurations.}
We adopt the three teacher--student pairs evaluated by TALAS \citep{talas}.
They comprise Qwen3-Embedding-4B $\rightarrow$ BERT-base, BGE-M3
$\rightarrow$ MiniLMv2-L6-H768, and Qwen3-Embedding-0.6B $\rightarrow$
MiniLMv2-L6-H384
\citep{devlin2019bert,wang2021minilmv2,chen2024m3,zhang2025qwen3embedding}.
These settings span students from 22M to 109M parameters and embedding
bottlenecks of $2560\rightarrow768$, $1024\rightarrow768$, and
$1024\rightarrow384$, respectively. The Qwen teachers use last-token pooling,
whereas BGE-M3 uses CLS pooling.

\paragraph{Distillation corpus.}
Following TALAS \citep{talas}, every method uses the same corpus of about 5K
unlabeled training texts from each of Emotion, WiC, and STS-B; no labels are
used. We call these three benchmarks source tasks and the other six non-source
tasks. ``Non-source'' denotes only absence from this corpus, not unrestricted
domain shift.

\paragraph{Baselines.}
The baselines span relational transfer (RKD), multi-stage endpoint distillation
(Stella), cross-tokenizer or dual-space transfer (CDM, DSKD), and hidden-state
alignment (EMO, TALAS)
\citep{park2019rkd,zhang2024jasper,chen2025cdm,zhang2024dskd,truong2025emo,talas}.
All share the corpus, model pairs, evaluator, five-epoch budget, batch size 128,
and learning-rate grid while retaining method-specific components. This is a
controlled-budget comparison, not an upper bound after per-pair tuning; full
implementation choices appear in Appendix~\ref{app:training-hyperparameters}.

\paragraph{Evaluation protocol.}
We evaluate frozen embeddings on classification (Banking77, Emotion, Tweet;
macro-F1), pair classification (MRPC, SciTail, WiC; average precision from
cosine similarity), and semantic similarity (SICK, STS12, STS-B; Spearman
correlation). Metrics are multiplied by 100 and averaged over three seeds.
Avg. is the unweighted mean of nine heterogeneous task metrics and is therefore
a descriptive index; Src and Non-src separately average the three source and
six non-source tasks.

\subsection{Main Results}
\label{sec:main-results}

\begin{table}[t]
    \caption{Nine-task results on the 15K distillation corpus. Distilled
    students report mean $\pm$ sample standard deviation over three seeds;
    best and second-best results are bold and underlined. GATE-KD uses
    $\lambda_{H_0}=0.5$.}
    \label{tab:main-results}
    \widetablebody
    \begin{tabular}{@{}lcccccccccccc@{}}
        \toprule
        & \multicolumn{3}{c}{Classification (F1)}
        & \multicolumn{3}{c}{Pair Classification (AP)}
        & \multicolumn{3}{c}{STS (Spearman)}
        & \multicolumn{3}{c}{Averages} \\
        \cmidrule(lr){2-4}\cmidrule(lr){5-7}\cmidrule(lr){8-10}\cmidrule(lr){11-13}
        Method & Banking77 & Tweet & Emotion & MRPC & SciTail & WiC
        & SICK & STS12 & STS-B & Src & Non-src & All \\
        \midrule
        \multicolumn{13}{@{}c}{\textit{(a) Qwen3-Embedding-4B $\rightarrow$ BERT-base (109M)}} \\
        \midrule
        Teacher
        & 93.26 & 71.04 & 69.08 & 84.80 & 91.77 & 66.79
        & 83.43 & 82.55 & 86.87 & 74.25 & 84.48 & 81.07 \\
        Student (pre-KD)
        & 84.33 & 67.08 & 47.79 & 77.36 & 67.74 & 59.22
        & 42.43 & 21.54 & 20.29 & 42.43 & 60.08 & 54.20 \\
        \midrule
        RKD
        & \meanstdstacked{89.76}{0.11} & \meanstdstacked{70.90}{0.31} & \meanstdstacked{59.98}{0.05}
        & \meanstdstacked{81.94}{0.26} & \meanstdstacked{75.59}{0.13} & \meanstdstacked{69.02}{0.18}
        & \meanstdstacked{72.93}{0.07} & \meanstdstacked{67.84}{0.18} & \meanstdstacked{73.72}{0.08}
        & \meanstdstacked{67.57}{0.05} & \meanstdstacked{76.49}{0.07} & \meanstdstacked{73.52}{0.06} \\
        Stella
        & \meanstdstacked{88.48}{0.30} & \meanstdstacked{68.94}{0.17} & \meanstdstacked{53.69}{0.50}
        & \meanstdstacked{83.95}{0.10} & \meanstdstacked{75.39}{0.54} & \meanstdstacked{\tablebest{71.93}}{0.09}
        & \meanstdstacked{69.87}{0.64} & \meanstdstacked{64.39}{0.21} & \meanstdstacked{70.86}{0.36}
        & \meanstdstacked{65.49}{0.17} & \meanstdstacked{75.17}{0.13} & \meanstdstacked{71.95}{0.06} \\
        CDM
        & \meanstdstacked{88.13}{0.23} & \meanstdstacked{69.03}{0.20} & \meanstdstacked{52.83}{0.50}
        & \meanstdstacked{83.00}{0.06} & \meanstdstacked{77.10}{0.27} & \meanstdstacked{71.84}{0.09}
        & \meanstdstacked{67.89}{0.37} & \meanstdstacked{61.06}{0.37} & \meanstdstacked{69.19}{0.42}
        & \meanstdstacked{64.62}{0.28} & \meanstdstacked{74.37}{0.12} & \meanstdstacked{71.12}{0.18} \\
        DSKD
        & \meanstdstacked{87.97}{0.16} & \meanstdstacked{68.90}{0.28} & \meanstdstacked{53.07}{0.31}
        & \meanstdstacked{83.03}{0.15} & \meanstdstacked{76.98}{0.59} & \meanstdstacked{\tablesecond{71.86}}{0.08}
        & \meanstdstacked{67.77}{0.66} & \meanstdstacked{61.11}{0.26} & \meanstdstacked{69.55}{0.12}
        & \meanstdstacked{64.83}{0.12} & \meanstdstacked{74.29}{0.14} & \meanstdstacked{71.14}{0.10} \\
        EMO
        & \meanstdstacked{86.89}{0.07} & \meanstdstacked{68.99}{0.11} & \meanstdstacked{52.02}{1.09}
        & \meanstdstacked{79.74}{0.12} & \meanstdstacked{76.37}{0.35} & \meanstdstacked{69.62}{0.16}
        & \meanstdstacked{67.89}{0.09} & \meanstdstacked{56.46}{0.47} & \meanstdstacked{64.54}{0.50}
        & \meanstdstacked{62.06}{0.49} & \meanstdstacked{72.72}{0.15} & \meanstdstacked{69.17}{0.24} \\
        TALAS
        & \meanstdstacked{\tablesecond{90.53}}{0.09} & \meanstdstacked{\tablesecond{73.77}}{0.21} & \meanstdstacked{\tablesecond{68.47}}{0.34}
        & \meanstdstacked{\tablebest{85.38}}{0.05} & \meanstdstacked{\tablebest{81.83}}{0.21} & \meanstdstacked{69.59}{0.14}
        & \meanstdstacked{\tablesecond{76.79}}{0.06} & \meanstdstacked{\tablesecond{70.79}}{0.19} & \meanstdstacked{\tablebest{78.50}}{0.08}
        & \meanstdstacked{\tablebest{72.19}}{0.07} & \meanstdstacked{\tablesecond{79.85}}{0.07} & \meanstdstacked{\tablesecond{77.29}}{0.04} \\
        \midrule
        \tablebest{GATE-KD}
        & \meanstdstacked{\tablebest{90.60}}{0.21} & \meanstdstacked{\tablebest{73.86}}{0.10} & \meanstdstacked{\tablebest{69.48}}{0.22}
        & \meanstdstacked{\tablesecond{85.31}}{0.11} & \meanstdstacked{\tablesecond{80.23}}{0.13} & \meanstdstacked{68.19}{0.12}
        & \meanstdstacked{\tablebest{77.80}}{0.14} & \meanstdstacked{\tablebest{75.15}}{0.07} & \meanstdstacked{\tablesecond{78.44}}{0.16}
        & \meanstdstacked{\tablesecond{72.04}}{0.08} & \meanstdstacked{\tablebest{80.49}}{0.03} & \meanstdstacked{\tablebest{77.67}}{0.03} \\
        \midrule
        \multicolumn{13}{@{}c}{\textit{(b) BGE-M3 $\rightarrow$ MiniLMv2-H768 (66M)}} \\
        \midrule
        Teacher
        & 93.49 & 73.75 & 69.03 & 85.80 & 91.87 & 61.46
        & 79.18 & 78.74 & 84.86 & 71.78 & 83.81 & 79.80 \\
        Student (pre-KD)
        & 81.18 & 71.30 & 48.85 & 79.66 & 65.06 & 60.89
        & 49.46 & 26.76 & 24.12 & 44.62 & 62.24 & 56.36 \\
        \midrule
        RKD
        & \meanstdstacked{89.14}{0.09} & \meanstdstacked{73.29}{0.09} & \meanstdstacked{\tablesecond{62.23}}{0.27}
        & \meanstdstacked{84.35}{0.06} & \meanstdstacked{78.74}{0.11} & \meanstdstacked{65.02}{0.10}
        & \meanstdstacked{73.20}{0.19} & \meanstdstacked{65.02}{0.29} & \meanstdstacked{71.97}{0.38}
        & \meanstdstacked{66.41}{0.07} & \meanstdstacked{77.29}{0.06} & \meanstdstacked{73.66}{0.04} \\
        Stella
        & \meanstdstacked{86.82}{0.28} & \meanstdstacked{70.06}{0.12} & \meanstdstacked{56.61}{0.74}
        & \meanstdstacked{83.02}{0.12} & \meanstdstacked{74.63}{0.12} & \meanstdstacked{\tablebest{69.72}}{0.05}
        & \meanstdstacked{69.08}{0.20} & \meanstdstacked{63.14}{0.28} & \meanstdstacked{67.10}{0.10}
        & \meanstdstacked{64.47}{0.27} & \meanstdstacked{74.46}{0.05} & \meanstdstacked{71.13}{0.11} \\
        CDM
        & \meanstdstacked{87.82}{0.19} & \meanstdstacked{69.99}{0.09} & \meanstdstacked{58.32}{0.51}
        & \meanstdstacked{83.44}{0.07} & \meanstdstacked{75.56}{0.40} & \meanstdstacked{\tablesecond{69.70}}{0.02}
        & \meanstdstacked{65.93}{0.16} & \meanstdstacked{60.84}{0.34} & \meanstdstacked{65.31}{0.15}
        & \meanstdstacked{64.45}{0.13} & \meanstdstacked{73.93}{0.06} & \meanstdstacked{70.77}{0.02} \\
        DSKD
        & \meanstdstacked{87.70}{0.16} & \meanstdstacked{70.35}{0.32} & \meanstdstacked{59.71}{0.61}
        & \meanstdstacked{82.99}{0.18} & \meanstdstacked{75.58}{0.12} & \meanstdstacked{69.59}{0.04}
        & \meanstdstacked{65.36}{0.37} & \meanstdstacked{61.21}{0.31} & \meanstdstacked{65.02}{0.18}
        & \meanstdstacked{64.77}{0.27} & \meanstdstacked{73.87}{0.10} & \meanstdstacked{70.83}{0.11} \\
        EMO
        & \meanstdstacked{87.54}{0.18} & \meanstdstacked{71.03}{0.14} & \meanstdstacked{58.52}{0.30}
        & \meanstdstacked{82.46}{0.18} & \meanstdstacked{74.73}{0.58} & \meanstdstacked{69.20}{0.22}
        & \meanstdstacked{64.71}{0.55} & \meanstdstacked{56.89}{0.27} & \meanstdstacked{61.25}{0.03}
        & \meanstdstacked{62.99}{0.19} & \meanstdstacked{72.89}{0.10} & \meanstdstacked{69.59}{0.12} \\
        TALAS
        & \meanstdstacked{\tablesecond{89.83}}{0.09} & \meanstdstacked{\tablebest{74.70}}{0.09} & \meanstdstacked{61.65}{0.23}
        & \meanstdstacked{\tablebest{86.62}}{0.05} & \meanstdstacked{\tablesecond{80.85}}{0.22} & \meanstdstacked{66.72}{0.07}
        & \meanstdstacked{\tablesecond{75.40}}{0.10} & \meanstdstacked{\tablesecond{67.08}}{0.31} & \meanstdstacked{\tablesecond{77.08}}{0.01}
        & \meanstdstacked{\tablesecond{68.48}}{0.07} & \meanstdstacked{\tablesecond{79.08}}{0.08} & \meanstdstacked{\tablesecond{75.55}}{0.06} \\
        \midrule
        \tablebest{GATE-KD}
        & \meanstdstacked{\tablebest{90.17}}{0.05} & \meanstdstacked{\tablesecond{73.74}}{0.24} & \meanstdstacked{\tablebest{64.83}}{0.72}
        & \meanstdstacked{\tablesecond{86.42}}{0.09} & \meanstdstacked{\tablebest{80.97}}{0.08} & \meanstdstacked{65.47}{0.09}
        & \meanstdstacked{\tablebest{76.77}}{0.03} & \meanstdstacked{\tablebest{68.27}}{0.04} & \meanstdstacked{\tablebest{77.63}}{0.12}
        & \meanstdstacked{\tablebest{69.31}}{0.21} & \meanstdstacked{\tablebest{79.39}}{0.06} & \meanstdstacked{\tablebest{76.03}}{0.11} \\
        \midrule
        \multicolumn{13}{@{}c}{\textit{(c) Qwen3-Embedding-0.6B $\rightarrow$ MiniLMv2-H384 (22M)}} \\
        \midrule
        Teacher
        & 92.51 & 71.28 & 65.80 & 83.38 & 90.09 & 66.84
        & 80.65 & 77.11 & 84.59 & 72.41 & 82.50 & 79.14 \\
        Student (pre-KD)
        & 67.03 & 69.45 & 42.90 & 78.39 & 64.55 & 59.58
        & 47.60 & 21.96 & 22.08 & 41.52 & 58.16 & 52.62 \\
        \midrule
        RKD
        & \meanstdstacked{\tablesecond{86.76}}{0.07} & \meanstdstacked{69.87}{0.21} & \meanstdstacked{56.31}{0.35}
        & \meanstdstacked{81.70}{0.03} & \meanstdstacked{76.03}{0.20} & \meanstdstacked{\tablesecond{68.56}}{0.22}
        & \meanstdstacked{69.21}{0.18} & \meanstdstacked{64.26}{0.19} & \meanstdstacked{70.25}{0.21}
        & \meanstdstacked{65.04}{0.20} & \meanstdstacked{74.64}{0.09} & \meanstdstacked{71.44}{0.10} \\
        Stella
        & \meanstdstacked{84.04}{0.04} & \meanstdstacked{68.73}{0.44} & \meanstdstacked{52.41}{0.28}
        & \meanstdstacked{81.50}{0.12} & \meanstdstacked{72.26}{0.32} & \meanstdstacked{\tablebest{68.65}}{0.16}
        & \meanstdstacked{64.64}{0.19} & \meanstdstacked{60.59}{0.12} & \meanstdstacked{64.27}{0.20}
        & \meanstdstacked{61.78}{0.13} & \meanstdstacked{71.96}{0.13} & \meanstdstacked{68.57}{0.10} \\
        CDM
        & \meanstdstacked{84.74}{0.21} & \meanstdstacked{69.52}{0.20} & \meanstdstacked{53.76}{0.95}
        & \meanstdstacked{81.15}{0.14} & \meanstdstacked{72.86}{0.30} & \meanstdstacked{68.29}{0.11}
        & \meanstdstacked{63.98}{0.36} & \meanstdstacked{60.43}{0.22} & \meanstdstacked{63.52}{0.24}
        & \meanstdstacked{61.86}{0.43} & \meanstdstacked{72.11}{0.19} & \meanstdstacked{68.69}{0.27} \\
        DSKD
        & \meanstdstacked{84.76}{0.10} & \meanstdstacked{69.56}{0.20} & \meanstdstacked{53.14}{0.29}
        & \meanstdstacked{81.15}{0.01} & \meanstdstacked{72.81}{0.09} & \meanstdstacked{68.33}{0.04}
        & \meanstdstacked{63.78}{0.06} & \meanstdstacked{60.24}{0.19} & \meanstdstacked{63.53}{0.18}
        & \meanstdstacked{61.67}{0.03} & \meanstdstacked{72.05}{0.06} & \meanstdstacked{68.59}{0.04} \\
        EMO
        & \meanstdstacked{83.54}{0.31} & \meanstdstacked{69.31}{0.06} & \meanstdstacked{54.48}{0.27}
        & \meanstdstacked{80.87}{0.13} & \meanstdstacked{72.00}{0.16} & \meanstdstacked{67.89}{0.09}
        & \meanstdstacked{62.19}{0.16} & \meanstdstacked{57.04}{0.17} & \meanstdstacked{61.24}{0.23}
        & \meanstdstacked{61.20}{0.09} & \meanstdstacked{70.82}{0.07} & \meanstdstacked{67.62}{0.02} \\
        TALAS
        & \meanstdstacked{86.39}{0.13} & \meanstdstacked{\tablesecond{71.81}}{0.11} & \meanstdstacked{\tablesecond{58.40}}{0.30}
        & \meanstdstacked{\tablebest{84.25}}{0.03} & \meanstdstacked{\tablesecond{79.29}}{0.11} & \meanstdstacked{66.82}{0.02}
        & \meanstdstacked{\tablesecond{70.42}}{0.10} & \meanstdstacked{\tablesecond{68.08}}{0.22} & \meanstdstacked{\tablesecond{73.34}}{0.10}
        & \meanstdstacked{\tablesecond{66.19}}{0.11} & \meanstdstacked{\tablesecond{76.71}}{0.04} & \meanstdstacked{\tablesecond{73.20}}{0.06} \\
        \midrule
        \tablebest{GATE-KD}
        & \meanstdstacked{\tablebest{87.77}}{0.07} & \meanstdstacked{\tablebest{72.30}}{0.19} & \meanstdstacked{\tablebest{63.45}}{0.43}
        & \meanstdstacked{\tablesecond{84.19}}{0.03} & \meanstdstacked{\tablebest{81.56}}{0.16} & \meanstdstacked{68.17}{0.05}
        & \meanstdstacked{\tablebest{72.59}}{0.03} & \meanstdstacked{\tablebest{68.91}}{0.03} & \meanstdstacked{\tablebest{74.79}}{0.05}
        & \meanstdstacked{\tablebest{68.80}}{0.16} & \meanstdstacked{\tablebest{77.89}}{0.01} & \meanstdstacked{\tablebest{74.86}}{0.06} \\
        \bottomrule
    \end{tabular}
\end{table}


GATE-KD has the highest numerical Avg. in all three configurations, exceeding
the next-best baseline by 0.38, 0.48, and 1.59 points. The first two are small
three-seed differences; the compact result is both larger and broader. Relative
to TALAS, GATE-KD is higher on 5/9, 6/9, and 8/9 task means, respectively
(Appendix~\ref{app:task-mean-differences}), and ranks first in 18 of 27
task--configuration comparisons overall. Pair classification remains mixed.

In configuration (c), endpoint-only GATE-KD reaches $74.26\pm0.02$, already
1.06 points above TALAS; $H_0$ adds 0.53 points. Thus the endpoint interface,
not the auxiliary, supplies most of the compact gain. The result also shows
that intermediate teacher states were not required in this comparison, without
implying that intermediate supervision is generally inferior.

On five corpus-disjoint BEIR retrieval benchmarks, GATE-KD ranks first on
average in (b) and (c) and second in (a), although all students remain below
their teachers (Appendix~\ref{app:retrieval-15k}). This is a transfer guardrail,
not a substitute for benchmark-disjoint general-text distillation.

\FloatBarrier
